Für das Grafikdesign haben wir verschiedene Bilder im Pixelartdesign erstellt, welche im Spielhintergrund zu sehen sind.
Der Spieler kann diese Hintergründe per Point-and-Click ändern, sodass sich die Spielfigur im Schulhaus bewegt. \par
Die Struktur der Hintergrundszenen im Spiel ist eine vereinfachte Version unserer Schule, es wurden hauptsächlich die Gänge und Räume eingebunden, die zu den 
spielrelevanten Fachschaften gehören. So wird im Spiel auf das B-Gebäude und den größten Teil des E-Gebäudes verzichtet. \\
Wir haben uns bei der Anordnung der Szenen für eine Ansicht vom großen Schulhof aus entschieden. Trotzdem startet die Spielfigur in der Foyerszene mit 
Sicht in die andere Richtung und der Raum des Hausmeisters ist nicht im Foyer sondern im Gang daneben. Solche Änderungen im Vergleich zur Realität folgen daraus,
dass das Spieldesign in 2D gegeben ist und man deshalb jeweils nur eine Seite eines Raumes darstellen kann. \\
Das allgemeine Aussehen der Gänge und Räume wurde im Spiel ebenfalls abgeändert. Da wir mit einem Pixelartdesign arbeiteten, haben wir uns bei der Gestaltung der 
Gänge und Räume auf charakteristische Darstellungen der verschiedenen Fachschaften reduziert. Auch die Farbgestaltung weicht von der Realität ab, für die 
Hintergründe nutzten wir grellere und vielfältigere Farben, um für eine erhöhten Freude am Spiel zu sorgen. \par
Das Verbinden der verschiedenen Szenen untereinander funktioniert so, dass es einen Szenenstapel gibt, in dem einzelne Szenen mithilfe von Objekten, auf die man 
klicken kann, ersetzt werden können. 
Für jede Szene wurde eine eigene Klasse erstellt, in der die jeweils benötigten Objekte eingefügt werden. \par
Eine weitere Aufgabe der Gruppe bestand darin, dafür zu sorgen, dass die Spielfigur in den Szenen auftaucht, wobei auf eine vielfältige Variabilität der Position
der Figur zu achten ist. \\ 
Die Spielfigur, welche als Objekt vorliegt, wird in jeder Szene individuell hinzugefügt.
Zur Umsetzung von verschiedenen Positionen der Spielfigur entwarfen wir eine Methode, welche bei jedem Szenenwechsel eine zufällige x-Koordinate für die 
Spielfigur erstellt. Die y-Koordinate ist fest definiert, da unsere Spielfigur nicht auf dem Bildschirm \glqq schweben\grqq\ soll.
